\subsection{变换与抽样}

主机上的 QR 分解将带种子的高斯矩阵分为正交因子和上三角因子，其中正交因子给出方形 $\Pi$。训练样本和数据库向量使用同一变换：
\begin{equation}
 Y_t=X_t\Pi^{\mathsf T},\qquad y=\Pi x,\qquad
 \|\Pi(q-x)\|_2^2=\|q-x\|_2^2
 \quad\text{when }\Pi^{\mathsf T}\Pi=I.
 \label{eq:4-transform}
\end{equation}
维度仍为 $d$。\cref{sec:4-6} 的验证区分存储矩阵的几何性质（K0）与通用矩阵乘法（general matrix multiplication，GEMM）的舍入。

主评测版本（v69） 存储 FP32 矩阵，但在 $d\geq1{,}536$ 时启用 cuBLAS TF32 tensor-core 运算。这影响两个 OpenAI 数据集在训练／add／查询中的旋转和质心 GEMM；实现提供的环境变量开关无法关闭这一模式。因此，FP32 存储并不表示完全采用 FP32 乘法。FP32 变体（v72） 修正了这一开关。敏感性实验复用同一个 TF32 训练缓存，改变 add／查询算术，并非以 FP32 重新训练模型（\cref{sec:7-precision}）。

主码／IVF 训练从数据库中等间隔选取样本，请求的样本数为 $n_t=39n_{\rm list}$：
\[
 n_t=\begin{cases}159{,}744,&\text{Vogue},\\319{,}488,&\text{arXiv/OpenAI},\\1{,}277{,}952,&\text{BGE/Stella}.\end{cases}
\]
这是 Faiss 的样本数警告阈值，而非量化质量保证 \cite{faiss_clustering}。

\subsection{解析主码本}

主码本由样本尺度解析生成，其 Cartesian 等级保留了编码器与查找表使用的坐标结构：
\begin{equation}
 \sigma^2=\frac{1}{n_t d}\sum_{i=0}^{n_t-1}\|x_i\|_2^2,
 \qquad D_s=d/M,\qquad K=2^B.
 \label{eq:4-scale}
\end{equation}
\[
 B=D_s=8\Longrightarrow
 \mathcal C_m=\{-a,+a\}^{8},\quad a=0.674\sigma.
\]
八个符号可放入一个 byte，支持 $O(d)$ 编码与拆分 LUT \cite{jhq}（\cref{sec:5-2}）。无约束 Lloyd 更新可能破坏这一 Cartesian 结构。对固定符号向量 $s\in\{-1,+1\}^d$，平方距离展开为
\begin{equation}
 \|q'-a s\|^2=\|q'\|^2+da^2-2a\langle q',s\rangle.
 \label{eq:4-sign-ranking}
\end{equation}
\[
 a>0\Longrightarrow
 \arg\min_s\|q'-as\|^2=\arg\max_s\langle q',s\rangle.
\]
因此，在 $a>0$ 时，主码排序由查询与符号向量的内积决定，无需数据服从高斯分布。残差仍依赖 $a$；零尺度和 FP32 并列需要另行处理。

\subsection{粗分配中的筛选与精确计算}

IVF 用间隔选取的旋转行初始化质心，并以 $s_\ell=\|\mu_\ell\|^2-2\langle y,\mu_\ell\rangle$ 运行八次迭代。支持的维度先用 int8 分数筛选候选，再在 FP32 中处理不确定情况。下面的引理说明区间筛选保留分配结果所需的条件。

\begin{resultblock}
\begin{lemma}[具有认证分数区间的分配]
\label{lem:4-assignment-bound}
设 $\widetilde y,\widetilde\mu_\ell$ 为反量化近似，$e=y-\widetilde y$，$f_\ell=\mu_\ell-\widetilde\mu_\ell$。保持质心范数精确时，近似分数
$\widetilde s_\ell=\|\mu_\ell\|^2-2\langle\widetilde y,\widetilde\mu_\ell\rangle$
的实数算术误差由下式界定：
\begin{equation}
 E_\ell=2\bigl(\|y\|\|f_\ell\|+\|e\|\|\widetilde\mu_\ell\|\bigr).
 \label{eq:4-assignment-bound}
\end{equation}
给定包含每个目标分数的有效区间 $[L_\ell,U_\ell]$，仅丢弃 $L_\ell>\min_j U_j$，并对其余分数精确求解，即可保留最小值及规定的并列规则。
\end{lemma}
\begin{proof}
使用 $e=y-\widetilde y$ 和 $f_\ell=\mu_\ell-\widetilde\mu_\ell$，
\[
 \begin{aligned}
 \langle y,\mu_\ell\rangle-
 \langle\widetilde y,\widetilde\mu_\ell\rangle
 &=\langle y,f_\ell\rangle
   +\langle e,\widetilde\mu_\ell\rangle,\\
 |s_\ell-\widetilde s_\ell|
 &\leq2\bigl(\|y\|\|f_\ell\|
       +\|e\|\|\widetilde\mu_\ell\|\bigr)=E_\ell .
 \end{aligned}
\]
对 $j^*\in\arg\min_j U_j$，
\[
 L_\ell>U_{j^*}
 \Longrightarrow
 s_\ell\geq L_\ell>U_{j^*}\geq s_{j^*}.
\]
严格不等号排除 $\ell$；等号情况仍可入选。
\end{proof}
\end{resultblock}

有效分数区间是关键条件：若无累加误差界，FP32 安全余量不能证明分配结果逐位一致。存档 FP16 尾部分配比较在 Vogue 的 0/14,838、BGE 的 6/31,748 和 Stella 的 3/16,359 个分配中报告差异。冷构建日志未标明实际采用的分配精度。

\subsection{残差拟合与迭代预算}

残差相对主码重建定义，即 $r=\Pi x-\widehat p(x)$。残差拟合使用 100,000 行。在每个子空间内，GPU 排序和 Lloyd 更新拟合 $K_r=2^{B_r}$ 个有序标量等级，在各坐标间共享，最多更新 2,000 次。

较早的十二单元 Vogue／OpenAI-3072 扫描（$B_r=4/8$，probe 为 32/128/256）固定主码／IVF 模型。在 25 次 Lloyd 迭代下，以 100K 行替代全部行最多损失 0.0018 召回率（十一个下降、一个持平）。中间样本量的结果并非单调，因此该证据在这一协议内支持有界残差训练。

历史 100K 训练从 25$\to$2,000 次更新增加了 0.0063--0.0076 召回率。尾部 bin 更新和 BGE／Stella 的结果仍在变化：达到上限并不证明收敛，也不违背精确 Lloyd 的下降性质。20,000 次更新时，历史 Stella 全量／100K 召回率为 0.9921/0.9920。

\subsection{训练工作量与列表不均衡}

训练工作量模型分别考虑变换生成、样本旋转、主码本生成和残差提取。以 $n_r$ 表示残差训练行数，变换部分为
\[
 W_{\rm QR}=O(d^3),\qquad W_{\rm rotate}=O(n_t d^2).
\]
直接粗分配每次迭代成本为 $O(n_t n_{\rm list}d)$，其余部分为
\[
 W_{\rm primary}=O(MKD_s),\qquad
 W_{\rm residual\ extraction}=O(n_rd),
\]
随后进行排序与标量更新。这些量表示操作数。同步诊断表明，BGE／Stella 训练约 83\% 的时间在 IVF k-means（约 2.97\,s）。该诊断使用一个 add 缓冲并强制同步，不同于 \cref{sec:7-7} 中正常的重叠测量。列表长度不均影响后续扫描工作量，对应 \cref{sec:5-5}；八次 IVF 迭代、间隔种子和有限抽样各自对不均衡的影响未单独测量。
